\documentclass[book.tex]{subfiles}
\begin{document}

\preface

These notes accompany a one-semester course on quantum optimal control taught at the Czech Technical University in Prague.
The course takes the point of view of a control engineer: a quantum device is a dynamical system with inputs (control fields), outputs (measurement statistics), and an internal state (a density operator) that is never observed directly.
The three questions of classical systems theory---what can be observed, what can be reached, and how much data is needed to learn the model---are asked again for closed and open quantum systems, and answered using the geometry of quantum states, channels, and control landscapes.
The final chapters turn the theory into algorithms for tomography and for the synthesis of control pulses.

\todo{Expand the preface once the course has been taught for the first time.}

\vspace{\baselineskip}
\begin{flushright}\noindent
Prague, \today\hfill {\it Jakub Mare{\v c}ek}
\end{flushright}

\end{document}
